% Generated by scripts/assemble.py; corrections belong in source/corrections.json.
\examhead{176}{西北大学 2024 年数学分析试卷}{166}
% p167-b01; source PDF page 167
\par\Needspace{5\baselineskip}
\blockanchor{p167-b01}
\majorquestion{一}(20 分) 求下列极限.\par
% p167-b02; source PDF page 167
\blockanchor{p167-b02}
(1) $\lim_{x\to 0}\left(\frac{2 + e^{\frac{1}{x}}}{1 + e^{\frac{4}{x}}} +\frac{\sin x}{|x|}\right).$\par
% p167-b03; source PDF page 167
\blockanchor{p167-b03}
(2) $\lim_{x\to 0}\frac{1}{x^3}\left[\left(\frac{2023 + \cos x}{2024}\right)^x -1\right].$\par
% p167-b04; source PDF page 167
\par\Needspace{5\baselineskip}
\par\Needspace{6\baselineskip}
\blockanchor{p167-b04}
\majorquestion{二}(20 分) 设 $x_{1} > 0$, 作迭代序列\par
\nopagebreak[4]
% p167-b05; source PDF page 167
\blockanchor{p167-b05}
\displayformula{x _ {n + 1} = \ln (1 + x _ {n}), n = 1, 2, \dots .}
% p167-b06; source PDF page 167
\blockanchor{p167-b06}
(1) 证明: $\lim_{n\to \infty}x_n$ 存在, 并求其值.\par
% p167-b07; source PDF page 167
\blockanchor{p167-b07}
(2) 求  $\lim_{n\to\infty}\frac{x_{n}x_{n+1}}{x_{n}-x_{n+1}}$ .\par
% p167-b08; source PDF page 167
\par\Needspace{5\baselineskip}
\blockanchor{p167-b08}
\majorquestion{三}(15 分) 设 $f(x)$ 在 $[-1, 1]$ 上有连续的二阶导数, 且 $f(0) = 0$.\par
% p167-b09; source PDF page 167
\blockanchor{p167-b09}
(1) 写出 $f(x)$ 的带拉格朗日型余项的一阶麦克劳林公式.\par
% p167-b10; source PDF page 167
\blockanchor{p167-b10}
(2) 证明: 至少存在一点 $\eta \in [-1, 1]$, 使得 $f''(\eta) = 3\int_{-1}^{1} f(x) \mathrm{d}x$.\par
% p167-b11; source PDF page 167
\par\Needspace{5\baselineskip}
\blockanchor{p167-b11}
\majorquestion{四}（15分）已知 $z = f(x,y)$ 的全微分为 $\mathrm{dz} = 2x\mathrm{dx} - 2y\mathrm{dy}$，且 $f(1,1) = 2$，求 $f(x,y)$ 在有界闭区域 $D: x^2 + \frac{y^2}{4} \leq 1$ 上的最大值与最小值.\par
% p167-b12; source PDF page 167
\par\Needspace{5\baselineskip}
\par\Needspace{6\baselineskip}
\blockanchor{p167-b12}
\majorquestion{五}(15 分) 设 $\Omega$ 为半球 $\{(x, y, z) \mid x^2 + y^2 + z^2 \leq 1, z \geq 0\}$, 计算\par
\nopagebreak[4]
% p167-b13; source PDF page 167
\blockanchor{p167-b13}
\displayformula{I = \iiint_ {\Omega} (\sqrt {2} x + z) ^ {2} \mathrm{d} x \mathrm{d} y \mathrm{d} z.}
% p167-b14; source PDF page 167
\par\Needspace{5\baselineskip}
\blockanchor{p167-b14}
\majorquestion{六}(15 分) 设 $P(x, y) = \frac{x(x^2 + y^2)^k}{y}$, $Q(x, y) = -\frac{x^2(x^2 + y^2)^k}{y^2}$.\par
% p167-b15; source PDF page 167
\blockanchor{p167-b15}
(1) 若积分 $\int_{L} P \, \mathrm{d}x + Q \, \mathrm{d}y$ 在 $D = \{(x, y) \mid y > 0\}$ 上与路径无关, 求 $k$ 的值.\par
% p167-b16; source PDF page 167
\par\Needspace{6\baselineskip}
\blockanchor{p167-b16}
(2) 对第 (1) 的 $k$, 在 $D$ 内求函数 $u(x, y)$, 使得 $\mathrm{d}u = P \mathrm{~d}x + Q \mathrm{~d}y$, 并计算 $I = \int_{(1,2)}^{(3,4)} P \mathrm{~d}x + Q \mathrm{~d}y$.
\majorquestion{七}(15 分) 计算第二类曲面积分\par
\nopagebreak[4]
% p167-b17; source PDF page 167
\blockanchor{p167-b17}
\displayformula{\iint_ {\Sigma} \mathrm{d} y \mathrm{d} z + z \mathrm{d} z \mathrm{d} x + \frac {e ^ {x}}{\sqrt {x ^ {2} + y ^ {2}}} \mathrm{d} x \mathrm{d} y.}
% p167-b18; source PDF page 167
\blockanchor{p167-b18}
其中 $\Sigma$ 是曲面 $z = \sqrt{x^2 + y^2}$ 介于两平面 $z = 1$ 和 $z = 2$ 之间的部分, 取下侧.\par
% p167-b19; source PDF page 167
\par\Needspace{5\baselineskip}
\blockanchor{p167-b19}
\majorquestion{八}(15 分) 求函数项级数 $\sum_{n=1}^{\infty}\left(x + \frac{1}{n}\right)^n$ 的收敛域, 并讨论和函数的连续性.\par
% p167-b20; source PDF page 167
\par\Needspace{5\baselineskip}
\par\Needspace{6\baselineskip}
\blockanchor{p167-b20}
\majorquestion{九}(20 分) 设 $\varphi(x) = \int_{0}^{+\infty} \frac{t}{2 + t^x} \, \mathrm{d}t (x \geq 3)$. 令\par
\nopagebreak[4]
% p167-b21; source PDF page 167
\blockanchor{p167-b21}
\displayformula{f _ {n} (x) = \sum_ {k = 0} ^ {n - 1} \frac {1}{n} \varphi \left(x + \frac {k}{n}\right), n = 1, 2, \dots .}
% p167-b22; source PDF page 167
\blockanchor{p167-b22}
证明: 函数列  $\{f_{n}(x)\}$  在  $[3, A]$  上一致收敛, 其中 A > 3.\par
